\section{Avance de Proyecto (Parte 02)}
\notaBorrador

\subsection{Propósito y estado general}
Este segundo avance continúa el documento del 25 de septiembre de 2026 y las
bitácoras de las semanas 1 y 2. Explica la recuperación del retraso inicial,
documenta la corrección de la \hu{03}, presenta las cuarenta historias con su
estado y actualiza el calendario hasta la entrega del 27 de noviembre.

El estado global indicado por el equipo es de \textbf{36 historias terminadas,
tres pendientes y una en proceso}. Se toma ese estado como base del corte
documental propuesto al 13 de noviembre. La asignación de cierres a semanas y
días es una reconstrucción para revisión; no se deduce del número de archivos,
pantallas o commits del repositorio.

\begin{center}
\begin{tabular}{lrrl}
\toprule Estado & Cantidad & Porcentaje & Historias \\\midrule
Terminadas & 36 & 90\,\% & Todas salvo 03, 10, 23 y 30 \\
En proceso & 1 & 2.5\,\% & \hu{30} \\
Pendientes & 3 & 7.5\,\% & \hu{03}, \hu{10}, \hu{23} \\\midrule
Total & 40 & 100\,\% & Sin duplicar identificadores \\\bottomrule
\end{tabular}
\end{center}

El 90\,\% expresa la proporción de historias terminadas. No representa el
porcentaje de esfuerzo consumido, la cobertura de pruebas ni la preparación
completa para producción. Las pruebas de aceptación y la entrega siguen
reservadas para el cierre del proyecto.

\subsection{Continuidad con las primeras dos semanas}
El calendario original programaba T1 del 28 de septiembre al 2 de octubre y
T2 del 5 al 12 de octubre. Sin embargo, la bitácora 1 registra del 6 al 9 de
octubre investigación de arquitectura, bases de datos y herramientas de
organización. La bitácora 2 registra del 12 al 16 de octubre la creación del
tablero, reasignación de actividades y pendientes de arquitectura; el 15 de
octubre se reconoce que T1 estaba atrasada.

Por tanto, no sería coherente afirmar que T1 terminó el 2 de octubre ni que
T2 concluyó el 12. La propuesta reprograma el cierre de T1 al 20 de octubre y
superpone el trabajo de cuentas, perfiles, catálogo y administración una vez
acordados los contratos mínimos de integración.

La recuperación se explica mediante tres medidas:
\begin{itemize}
\item Cerrar decisiones de arquitectura y datos antes de ampliar los módulos,
para reducir cambios repetidos de formularios y contratos.
\item Distribuir frentes entre los ocho integrantes conservando al responsable
de cada historia y permitiendo apoyo cruzado.
\item Adelantar exploración e interacción del consumidor mediante entregas
parciales integrables y revisión diaria de bloqueos.
\end{itemize}
No se atribuye la recuperación a horas extraordinarias no registradas ni se
considera terminado un módulo únicamente por disponer de un prototipo.

\subsection{Corrección de la historia HU-03}
\textbf{Redacción anterior:} como emprendedor, publicar una foto del menú del
negocio para mostrar de un vistazo las ofertas.

\textbf{Redacción corregida, confirmada para este avance:} como cliente,
quiero acceder al sitio web sin registrarme, para consultar públicamente los
emprendimientos y sus productos.

La nueva historia corresponde al rol cliente o consumidor. No agrega cuentas
de clientes, contraseñas de consumidores ni un nuevo inicio de sesión.
El acceso del emprendedor se conserva en \hu{40}. La \hu{03} se retira del
grupo T4 (catálogo) y se incorpora al grupo T2 como acceso público; su
aceptación se coordina con T5 y T6, que contienen las pantallas y recorridos
a consultar.

\textbf{Criterios de aceptación propuestos:}
\begin{enumerate}[nosep,leftmargin=*]
\item Abrir la página principal y el directorio sin registro, token o sesión.
\item Abrir directamente el perfil de un negocio y consultar catálogo,
precios, horarios y medios de contacto públicos.
\item Usar búsqueda y categorías sin que se solicite una cuenta; las
funciones de zona se aceptarán cuando cierren \hu{10} y \hu{23}.
\item Mantener protegidas las rutas de administración y del emprendedor.
\item Respetar la privacidad de ubicación: si el negocio la oculta, no
exponer zona, identificadores ni referencias de localización.
\end{enumerate}

\textbf{Estado: pendiente.} Se revisó el alcance y se prepararon escenarios,
pero no se declara aceptado el recorrido completo sin sesión. Esto permite
que \hu{14} y \hu{18} estén terminadas como funciones de listado y perfil
sin declarar cerrada automáticamente la aceptación transversal de \hu{03}.

Esta corrección se registra en el entregable académico solicitado. La
especificación técnica actual del repositorio todavía describe la foto de
menú como \hu{03}; este documento no acredita que esa especificación o su
matriz técnica ya hayan sido migradas. Su conciliación queda incluida en el
cierre documental previo a la entrega.

\subsection{Criterio de estados y numeración}
\begin{itemize}
\item \textbf{Terminada:} incluida en las 36 que el equipo reporta como
realizadas; debe asociarse su evidencia antes de aprobar este borrador.
\item \textbf{En proceso:} tiene desarrollo parcial, pero falta un criterio
necesario para su cierre. En \hu{30} queda pendiente la cercanía.
\item \textbf{Pendiente:} todavía no se acepta la funcionalidad completa,
aunque se haya analizado, diseñado o preparado su integración.
\end{itemize}

Se conservan los identificadores 01--40 del primer avance. En su tabla de
prioridades, \hu{20} corresponde a moderación; en la asignación de responsables
aparece también revisión de solicitudes. Para evitar duplicación, este avance
mantiene \hu{20} como moderación, con Ángel, y concentra revisión y aprobación
de solicitudes en \hu{21}, con Denzel y apoyo de Ángel. Los desgloses técnicos
20A y 20B no se suman como historias adicionales.

\clearpage
\subsection{Matriz de las cuarenta historias de usuario}
Se conservan los responsables del primer avance. La prioridad y complejidad
de las historias sin cambio siguen su tabla original. Para \hu{03} corregida
se propone prioridad alta y complejidad media, sujetas a revisión del equipo.
En la última columna, S3--S6 indica la semana de cierre propuesta.

\begingroup\footnotesize
\begin{longtable}{@{}L{1.05cm}L{7.65cm}L{1.5cm}L{1.5cm}L{2.0cm}L{1.2cm}@{}}
\rowcolor{dorado}
\textcolor{marino}{\textbf{ID}}&
\textcolor{marino}{\textbf{Como / quiero / para}}&
\textcolor{marino}{\textbf{Resp.}}&
\textcolor{marino}{\textbf{Prior. / complej.}}&
\textcolor{marino}{\textbf{Estado}}&
\textcolor{marino}{\textbf{Cierre}}\\\midrule\endfirsthead
\rowcolor{suave}\textbf{ID}&\textbf{Como / quiero / para}&\textbf{Resp.}&
\textbf{Prior. / complej.}&\textbf{Estado}&\textbf{Cierre}\\\midrule\endhead
\bottomrule\endfoot
\hu{01} & \textbf{Emprendedor:} Solicitar una cuenta para registrar mi negocio y enviar mis datos para revisión. & Fernanda & Alta / Alta & \terminada & S3 \\\addlinespace[4pt]
\hu{02} & \textbf{Emprendedor:} Publicar el catálogo con imagen y descripción para mostrar a los clientes las ofertas de mis productos. & Fernanda & Alta / Alta & \terminada & S4 \\\addlinespace[4pt]
\hu{03} & \textbf{Cliente:} Acceder al sitio web sin registrarme para consultar públicamente los emprendimientos y sus productos. & Fernanda & Alta / Media & \pendiente & --- \\\addlinespace[4pt]
\hu{04} & \textbf{Emprendedor:} Agregar una foto de perfil para identificar mi negocio mediante una foto o logotipo. & Fernanda & Baja / Baja & \terminada & S3 \\\addlinespace[4pt]
\hu{05} & \textbf{Emprendedor:} Clasificar mi negocio para facilitar su localización por categoría. & Fernanda & Baja / Baja & \terminada & S3 \\\addlinespace[4pt]
\hu{06} & \textbf{Emprendedor:} Agregar una descripción del negocio para explicar los productos o servicios que ofrezco. & Orlando & Alta / Baja & \terminada & S3 \\\addlinespace[4pt]
\hu{07} & \textbf{Emprendedor:} Indicar el precio de cada producto para informar su costo antes del contacto. & Orlando & Alta / Baja & \terminada & S4 \\\addlinespace[4pt]
\hu{08} & \textbf{Emprendedor:} Editar productos del catálogo para mantener actualizada mi oferta. & Orlando & Alta / Baja & \terminada & S4 \\\addlinespace[4pt]
\hu{09} & \textbf{Emprendedor:} Eliminar productos del catálogo para retirar productos que ya no ofrezco. & Orlando & Alta / Baja & \terminada & S4 \\\addlinespace[4pt]
\hu{10} & \textbf{Emprendedor:} Publicar una ubicación aproximada para facilitar que me encuentren dentro del campus. & Orlando & Media / Media & \pendiente & --- \\\addlinespace[4pt]
\hu{11} & \textbf{Emprendedor:} Pausar productos para indicar su falta de disponibilidad temporal. & Mario & Media / Baja & \terminada & S4 \\\addlinespace[4pt]
\hu{12} & \textbf{Emprendedor:} Desactivar la ubicación pública para proteger mi privacidad cuando no deseo mostrarla. & Mario & Media / Baja & \terminada & S4 \\\addlinespace[4pt]
\hu{13} & \textbf{Emprendedor:} Agregar medios de contacto para permitir que los clientes se comuniquen conmigo. & Mario & Alta / Baja & \terminada & S3 \\\addlinespace[4pt]
\hu{14} & \textbf{Consumidor:} Visualizar los emprendimientos disponibles para conocer los negocios de la universidad. & Mario & Alta / Media & \terminada & S5 \\\addlinespace[4pt]
\hu{15} & \textbf{Consumidor:} Buscar un emprendimiento por nombre para encontrar rápidamente un negocio específico. & Mario & Alta / Media & \terminada & S5 \\\addlinespace[4pt]
\hu{16} & \textbf{Administrador:} Gestionar las categorías del sistema para mantener organizada la clasificación de los negocios. & Denzel & Media / Media & \terminada & S4 \\\addlinespace[4pt]
\hu{17} & \textbf{Emprendedor:} Indicar si mi negocio está abierto o cerrado para informar si puedo atender en ese momento. & Ángel & Alta / Baja & \terminada & S3 \\\addlinespace[4pt]
\hu{18} & \textbf{Consumidor:} Consultar el perfil de un emprendimiento para conocer sus productos, precios, horarios y contactos. & Ángel & Alta / Media & \terminada & S5 \\\addlinespace[4pt]
\hu{19} & \textbf{Consumidor:} Filtrar emprendimientos por categoría para encontrar opciones acordes con lo que busco. & Ángel & Alta / Media & \terminada & S5 \\\addlinespace[4pt]
\hu{20} & \textbf{Administrador:} Ocultar publicaciones que incumplan las reglas para evitar la exposición de contenido inapropiado o engañoso. & Ángel & Alta / Media & \terminada & S4 \\\addlinespace[4pt]
\hu{21} & \textbf{Administrador:} Revisar y resolver solicitudes de registro para aprobar o rechazar cuentas verificando pertenencia a la UNISON. & Denzel & Alta / Alta & \terminada & S4 \\\addlinespace[4pt]
\hu{22} & \textbf{Administrador:} Suspender temporalmente una cuenta para aplicar sanciones por incumplimiento de las reglas. & Denzel & Media / Alta & \terminada & S4 \\\addlinespace[4pt]
\hu{23} & \textbf{Consumidor:} Buscar emprendimientos por zona del campus para localizar negocios cercanos sin requerir una ubicación exacta. & Ángel & Media / Media & \pendiente & --- \\\addlinespace[4pt]
\hu{24} & \textbf{Administrador:} Gestionar los reportes de consumidores para tomar decisiones sobre cuentas denunciadas. & Denzel & Media / Alta & \terminada & S6 \\\addlinespace[4pt]
\hu{25} & \textbf{Administrador:} Dar mantenimiento general a la plataforma para asegurar su correcto funcionamiento. & Denzel & Media / Media & \terminada & S6 \\\addlinespace[4pt]
\hu{26} & \textbf{Consumidor:} Reportar un negocio con información falsa o inapropiada para contribuir a un directorio confiable. & Leinad & Media / Baja & \terminada & S5 \\\addlinespace[4pt]
\hu{27} & \textbf{Emprendedor:} Recibir notificación de aprobación o rechazo para conocer el resultado de mi solicitud. & Leinad & Alta / Baja & \terminada & S4 \\\addlinespace[4pt]
\hu{28} & \textbf{Emprendedor:} Eliminar mi cuenta para dejar de aparecer en el directorio al cerrar mi negocio. & Leinad & Media / Baja & \terminada & S5 \\\addlinespace[4pt]
\hu{29} & \textbf{Consumidor:} Contactar al emprendedor desde su perfil para comunicarme directamente por teléfono o redes sociales. & Leinad & Alta / Media & \terminada & S5 \\\addlinespace[4pt]
\hu{30} & \textbf{Consumidor:} Ordenar resultados por nombre, cercanía o estado para encontrar más rápido el emprendimiento que me interesa. & Leinad & Media / Alta & \enproceso & --- \\\addlinespace[4pt]
\hu{31} & \textbf{Consumidor:} Guardar emprendimientos como favoritos para volver a encontrar los negocios que me interesaron. & Paulina & Media / Media & \terminada & S5 \\\addlinespace[4pt]
\hu{32} & \textbf{Emprendedor:} Recuperar mi contraseña para volver a acceder sin crear otra cuenta. & Paulina & Alta / Media & \terminada & S4 \\\addlinespace[4pt]
\hu{33} & \textbf{Consumidor:} Filtrar productos por rango de precio para encontrar opciones acordes con mi presupuesto. & Paulina & Media / Media & \terminada & S5 \\\addlinespace[4pt]
\hu{34} & \textbf{Emprendedor:} Previsualizar mi perfil como consumidor para revisar la información antes de promocionarlo. & Paulina & Media / Baja & \terminada & S6 \\\addlinespace[4pt]
\hu{35} & \textbf{Administrador:} Consultar los reportes de consumidores para revisar negocios reportados y decidir acciones. & Paulina & Alta / Media & \terminada & S6 \\\addlinespace[4pt]
\hu{36} & \textbf{Emprendedor:} Establecer el horario de atención para informar cuándo pueden encontrarme. & Marco & Media / Baja & \terminada & S3 \\\addlinespace[4pt]
\hu{37} & \textbf{Consumidor:} Buscar un producto por su nombre para localizar los emprendimientos que lo ofrecen. & Marco & Media / Media & \terminada & S5 \\\addlinespace[4pt]
\hu{38} & \textbf{Consumidor:} Filtrar negocios por horario de atención para encontrar opciones disponibles a la hora que necesito. & Marco & Media / Media & \terminada & S6 \\\addlinespace[4pt]
\hu{39} & \textbf{Consumidor:} Filtrar negocios por estado abierto o cerrado para evitar contactar negocios que no están atendiendo. & Marco & Alta / Media & \terminada & S6 \\\addlinespace[4pt]
\hu{40} & \textbf{Emprendedor:} Iniciar sesión en mi cuenta para administrar de forma segura la información de mi negocio. & Marco & Alta / Media & \terminada & S3 \\\addlinespace[4pt]
\end{longtable}\endgroup

\subsection{Resumen por responsable}
Los apoyos de integración no cambian al responsable principal. La cantidad de
historias es una medida de distribución; no implica igual carga, pues existen
diferencias de complejidad y dependencia.

\begin{center}\small
\begin{tabular}{lrrrr}
\toprule Responsable & Terminadas & En proceso & Pendientes & Total \\\midrule
Fernanda & 4 & 0 & 1 & 5 \\
Orlando & 4 & 0 & 1 & 5 \\
Mario & 5 & 0 & 0 & 5 \\
Ángel & 4 & 0 & 1 & 5 \\
Denzel & 5 & 0 & 0 & 5 \\
Leinad & 4 & 1 & 0 & 5 \\
Paulina & 5 & 0 & 0 & 5 \\
Marco & 5 & 0 & 0 & 5 \\\midrule
Total & 36 & 1 & 3 & 40 \\\bottomrule
\end{tabular}
\end{center}

\subsection{Evolución propuesta entre las semanas 3 y 6}
No se asignan cierres funcionales a las primeras dos semanas porque los
registros aportados no los documentan. Esto no elimina su trabajo de
investigación, organización y diseño.

\begin{center}\small
\begin{tabularx}{\linewidth}{@{}lL{2.65cm}rrY@{}}
\toprule Semana & Periodo & Nuevas & Acumulado & Resultado principal \\\midrule
3 & 19--23 oct. & 8 & 8 (20\,\%) & Arquitectura, cuentas y perfil básico. \\
4 & 26--30 oct. & 12 & 20 (50\,\%) & Catálogo, cuentas restantes y administración inicial. \\
5 & 2--6 nov. & 10 & 30 (75\,\%) & Directorio, búsqueda e interacción. \\
6 & 9--13 nov. & 6 & 36 (90\,\%) & Reportes, mantenimiento, vista previa y filtros. \\\bottomrule
\end{tabularx}
\end{center}

La suma 8 + 12 + 10 + 6 = 36 es una propuesta de distribución del avance
global reportado. Las bitácoras detallan la participación diaria. Todas las
historias aparecen en el seguimiento: las pendientes se trabajan mediante
análisis y preparación, y no se incorporan a la suma de cierres.

\clearpage
\subsection{Cronograma actualizado de recuperación}
\subsubsection{Calendario y dependencias revisadas}
Se mantiene la ventana global del 28 de septiembre al 27 de noviembre:
61 días naturales y 45 jornadas de lunes a viernes, con conteo inclusivo y sin
descontar feriados. La numeración de bitácoras continúa los registros reales:
semana 1, 5--9 de octubre (con entradas desde el día 6); semana 2, 12--16;
semana 3, 19--23; semana 4, 26--30; semana 5, 2--6 de noviembre; semana 6,
9--13. La semana 7 corresponde al 16--20 y la semana 8 al 23--27.

El periodo del 28 de septiembre al 2 de octubre pertenece a la línea base,
pero no se presenta como una semana de ejecución validada. Las fechas nuevas
son ventanas de trabajo y aceptación, no estimaciones de esfuerzo individual.

Se reemplaza la dependencia de módulos completamente terminados por
entregables mínimos acordados: un catálogo puede avanzar con un contrato de
perfil estable, mientras la vista previa continúa en integración. La búsqueda
general puede avanzar sin zonas; el filtro por zona y la cercanía siguen
condicionados a ubicación pública. Así se explica el trabajo paralelo sin
afirmar que desaparecieron las dependencias técnicas.

\begingroup\scriptsize
\begin{longtable}{@{}L{0.9cm}L{2.4cm}L{1.8cm}L{2.0cm}L{0.7cm}L{3.3cm}L{4.5cm}@{}}
\rowcolor{dorado}\textbf{Tarea}&\textbf{Alcance}&\textbf{Original}&
\textbf{Ajustado}&\textbf{Días}&\textbf{Coordinación}&\textbf{Estado al corte / condición}\\\midrule\endfirsthead
\rowcolor{suave}\textbf{Tarea}&\textbf{Alcance}&\textbf{Original}&
\textbf{Ajustado}&\textbf{Días}&\textbf{Coordinación}&\textbf{Estado al corte / condición}\\\midrule\endhead
\bottomrule\endfoot
T1 & Arquitectura y BD & 28 sep.--2 oct. & 6--20 oct. & 11 &
Leinad, Denzel & Cierre propuesto el 20 oct.; contratos mínimos habilitan integración. \\
T2 & Cuentas y acceso público & 5--12 oct. & 19 oct.--17 nov. & 22 &
Mario, Fernanda; Marco y Paulina & Cuentas concluidas en S4; HU-03 corregida pendiente, prevista 16--17 nov. \\
T3 & Perfil del emprendimiento & 13--22 oct. & 20 oct.--12 nov. & 18 &
Ángel, Orlando; Mario, Paulina y Marco & Terminada al cerrar vista previa en S6. Integra contactos y horarios. \\
T4 & Catálogo y productos & 23--30 oct. & 23--30 oct. & 6 &
Marco, Paulina; Fernanda, Orlando y Mario & Terminada en S4. Incluye HU-02, 07, 08, 09 y 11; se retira HU-03. \\
T5 & Exploración, búsqueda y filtros & 2--16 nov. & 29 oct.--20 nov. & 17 &
Leinad, Denzel; Mario, Ángel, Paulina y Marco & En proceso por HU-30; los filtros de horario y estado cierran en S6. \\
T6 & Interacción del consumidor & 17--20 nov. & 2--5 nov. & 4 &
Marco, Paulina; Leinad & Terminada en S5: reporte, contacto y favoritos. \\
T7 & Administración & 13--22 oct. & 21 oct.--12 nov. & 17 &
Fernanda, Ángel; Denzel y Paulina & Terminada en S6 con reportes y mantenimiento. \\
T8 & Ubicación & 23--30 oct. & 22 oct.--20 nov. & 22 &
Orlando, Mario; Ángel & HU-12 terminada; HU-10 y HU-23 pendientes. Requiere integrar publicación antes de búsqueda. \\
T9 & Pruebas y entrega & 23--27 nov. & 9--27 nov. & 15 &
Denzel, Fernanda; todo el equipo & Revisión incremental 9--20 nov.; aceptación final y entrega 23--27 nov. \\
\end{longtable}\endgroup

T2 amplía explícitamente su alcance por el cambio de \hu{03}. Los contratos de
cuentas necesarios para T3 y T7 se habilitan en octubre; el cierre de acceso
público no bloquea el desarrollo de esos módulos. Las coordinaciones por
tarea se conservan del primer avance, con apoyo de los responsables
individuales que aparecen en la matriz.

Las duraciones anteriores se cuentan como lunes a viernes dentro de cada
ventana; no se suman para calcular la duración del proyecto porque varias
tareas se superponen. La línea base fechada se conserva como referencia;
no se mezclan sus días con las estimaciones originales de 10, 15 o 7 días
que no coincidían con el Gantt.

\subsubsection{Vista semanal del cronograma actualizado}
El sombreado azul indica ventana de trabajo o revisión. La columna inicial
resume las dos primeras semanas registradas. La fecha final permanece el
27 de noviembre.

\begin{center}
\begin{tikzpicture}[x=1.65cm,y=0.68cm,font=\scriptsize]
\fill[suave] (-2.8,0.45) rectangle (7.5,-9.5);
\foreach \x/\label in {0/S1--S2,1/S3,2/S4,3/S5,4/S6,5/S7,6/S8}{
 \node[font=\scriptsize\bfseries] at (\x,0) {\label};
}
\foreach \y/\label in {1/T1: arquitectura,2/T2: cuentas y acceso,3/T3: perfil,4/T4: catálogo,5/T5: búsqueda,6/T6: interacción,7/T7: administración,8/T8: ubicación,9/T9: pruebas y entrega}{
 \node[anchor=east] at (-0.65,-\y) {\label};
}
\foreach \row/\start/\finish in {1/0/1,2/1/5,3/1/4,4/1/2,5/2/5,6/3/3,7/1/4,8/1/5,9/4/6}{
 \fill[marino] (\start-0.4,-\row-0.22) rectangle (\finish+0.4,-\row+0.22);
}
\end{tikzpicture}
\end{center}
{\small S3: 19--23 oct.; S4: 26--30 oct.; S5: 2--6 nov.; S6: 9--13 nov.;
S7: 16--20 nov.; S8: 23--27 nov. El gráfico representa ventanas semanales;
las fechas exactas se especifican en la tabla.}

\subsubsection{Dependencias que condicionan el cierre}
La cadena funcional que requiere mayor atención es:
\[
\hu{10}\ \longrightarrow\ \hu{23}\ \longrightarrow\
\text{cercanía de \hu{30}}\ \longrightarrow\ \text{aceptación final}.
\]
La \hu{03} corregida se cierra en paralelo mediante la revisión de acceso
público. Ocultar la ubicación (\hu{12}) ya está terminada: protege los datos
privados y no equivale a haber terminado la captura, publicación o búsqueda
por zona.

No se reutilizan las holguras de 21 o 30 días del plan anterior. Al modificar
las dependencias, el alcance de T2 y las ventanas de trabajo, esas holguras y
las rutas críticas de la línea base dejan de describir el plan actualizado.
La cadena anterior identifica una dependencia de cierre; no es un cálculo
CPM con duraciones optimistas o un nuevo valor de holgura.

\subsection{Plan de cierre de historias abiertas}
\begingroup\small
\begin{longtable}{@{}L{1.5cm}L{1.6cm}L{2.2cm}L{4.2cm}L{6.5cm}@{}}
\rowcolor{dorado}\textbf{Historia}&\textbf{Resp.}&\textbf{Ventana}&
\textbf{Trabajo restante}&\textbf{Condición para terminar}\\\midrule\endfirsthead
\rowcolor{suave}\textbf{Historia}&\textbf{Resp.}&\textbf{Ventana}&
\textbf{Trabajo restante}&\textbf{Condición para terminar}\\\midrule\endhead
\bottomrule\endfoot
\hu{03}&Fernanda&16--17 nov.&
Completar aceptación de acceso público sin cuenta.&
Recorrido de inicio, directorio y perfil sin sesión; rutas de gestión
protegidas. Los casos por zona se revisan tras HU-23.\\
\hu{10}&Orlando&16--18 nov.&
Integrar captura y publicación de ubicación aproximada.&
Guardar zona y referencia, mostrar únicamente datos autorizados y conservar
la privacidad al desactivar visibilidad.\\
\hu{23}&Ángel&19--20 nov.&
Integrar búsqueda por zona con datos publicados.&
Resultados de la zona seleccionada; ausencia de negocios con ubicación
privada. Depende de la integración de HU-10.\\
\hu{30}&Leinad&19--20 nov.&
Completar cercanía e integrar todas las opciones de orden.&
Orden alfabético, por estado y por cercanía coherentes; negocios sin
ubicación pública sin exposición de datos. El cierre depende de HU-10 y HU-23.\\
\end{longtable}\endgroup

La implementación de cercanía puede avanzar con el contrato de zonas desde
el 19 de noviembre; su aceptación final se realiza después de validar
\hu{23}, el 20. Las cuatro ventanas son metas, no cierres consumados. Si no se
completa un criterio, la historia conserva su estado abierto y se ajusta el
plan de aceptación sin cambiar retrospectivamente las bitácoras.

\subsection{Riesgos y controles de seguimiento}
\begin{itemize}
\item \textbf{Ubicación y cercanía:} si HU-10 se retrasa, afecta HU-23 y el
cierre de HU-30. Orlando, Ángel y Leinad revisan diariamente datos y bloqueos.
\item \textbf{Cambio de alcance de HU-03:} evitar confundir acceso público con
autenticación de clientes. Fernanda registra criterios, casos sin sesión y
la conciliación de la documentación técnica.
\item \textbf{Integración paralela:} coordinar contratos y revisar flujos
completos para detectar incompatibilidades antes de las pruebas finales.
\item \textbf{Diferencias entre borrador y ejecución:} confirmar cada entrada,
fecha y cierre con el responsable; asociar evidencia antes de aprobar.
\end{itemize}

\subsection{Aceptación final y entrega}
Del 23 al 27 de noviembre se reserva la aceptación del conjunto:
\begin{itemize}
\item Revisar los recorridos de cliente sin registro, emprendedor y
administrador, incluyendo rutas protegidas y estados de cuenta.
\item Verificar privacidad de ubicación y credenciales, validaciones,
mensajes claros y consistencia de horarios.
\item Revisar uso móvil, navegación a contacto y correspondencia entre
perfil público y vista previa.
\item Ejecutar las verificaciones técnicas del proyecto y registrar
resultados e incidencias, sin anticipar que pasarán.
\item Conciliar HU-03 y la numeración de HU-20 en los documentos de cierre,
preparar evidencia y consolidar la entrega del 27 de noviembre.
\end{itemize}

El objetivo de la recuperación es llegar al cierre de las cuarenta historias
antes de la aceptación final. Este avance conserva el corte de 36 terminadas
y cuatro abiertas; no presenta el objetivo futuro como resultado ya obtenido.

\subsection{Fuentes de elaboración y revisión pendiente}
Se consultaron los siguientes materiales proporcionados por el equipo:
\begin{enumerate}[nosep,leftmargin=*]
\item \textit{Copia de AvanceDeProyecto (Parte 01)}, fechado el 25 de
septiembre de 2026: historias, prioridades, complejidad y responsables.
\item \textit{Bitácora semana 1}: entradas del 6 al 9 de octubre; portada
fechada el 22 de octubre.
\item \textit{Bitácora semana 2}: entradas del 12 al 16 de octubre; portada
fechada el 22 de octubre.
\item Descripción del cronograma original: calendario del 28 de septiembre
al 27 de noviembre y advertencia sobre diferencias con las duraciones
estimadas.
\item Confirmaciones del equipo para este documento: 36 terminadas,
HU-03, HU-10 y HU-23 pendientes, HU-30 en proceso; nuevo alcance de HU-03
como consulta sin registro y semana 6 completa como borrador.
\end{enumerate}

Se consultaron también la especificación técnica y el plan de construcción
del repositorio como referencia de nombres y dependencias. El reporte de
avance no sustituye la verificación de aceptación de cada funcionalidad.

Antes de convertir este documento en entrega definitiva, cada integrante
debe confirmar las actividades atribuidas, la semana de cierre y la evidencia
correspondiente. La portada y las fechas representan el corte propuesto del
13 de noviembre, no una aprobación ya firmada.
